\section{作品概述}
\subsection{背景介绍}
农业是我国第一用水大户，灌溉用水效率的提升是节水型社会建设的关键环节。华北平原集中了冬小麦—夏玉米一年两熟的主要产区，降水年内分配不均，春季与初夏的作物需水高峰往往与少雨期重叠，灌溉高度依赖地下水，长期超采已形成大范围地下水漏斗。在这一背景下，“什么时候灌、灌多少”不再只是农事经验问题，而是同时关系粮食稳产与水资源可持续利用的管理问题。

国际上通行的灌溉决策框架以FAO-56为代表：先由气象条件计算参考蒸散ET\textsubscript{0}，再乘以作物系数得到作物需水，结合根区土壤水量平衡判断是否需要补水\cite{fao,pereira}。这一框架机理清晰，但在基层落地时面临两点困难：其一，决策需要的是\emph{次日}乃至未来几天的需水，而不是已经发生的蒸散；其二，任何预测都有误差，单一预测值无法告诉使用者“如果预测偏低会怎样”。经验灌溉往往通过多灌来对冲不确定性，造成水资源浪费和深层渗漏；按平均值精确灌溉，又可能在需水突增时出现缺水胁迫\cite{fereres}。

近年来，机器学习方法在日参考蒸散估算与预测上取得了较多进展。随机森林、梯度提升树等集成方法在中国不同气候区的ET\textsubscript{0}建模中表现稳定，在气象输入不完整时也能保持较好精度\cite{fan,feng}。但已有工作大多止步于“预测精度”，较少回答两个与灌溉直接相关的问题：预测误差有多大、能否被量化；预测结果进入灌溉规则后，用水量与缺水风险如何变化。这正是本作品切入的位置。

\subsection{作品名称释义}
作品命名为“墒策”。“墒”（shāng）是我国传统农业术语，指土壤中适宜作物生长的水分状况，生产上常说“看墒情”“保墒”“抢墒播种”，各级农业部门也长期开展土壤墒情监测。“策”指决策。墒策即“依墒情定灌溉之策”：系统以根区土壤水分亏缺为核心状态量，把次日需水预测及其误差区间转化为具体的补水建议。副标题“可信需水预测与灌溉决策”中的“可信”，指预测结果附带经独立年份校准的误差区间，并且每一条建议都能追溯到输入数据、模型版本和计算过程。

\subsection{目标定位}
本作品的目标是构建一个面向农技人员与农业信息化教学的灌溉决策工具，具体包括：
\begin{enumerate}
\item \textbf{次日需水预测}：仅使用目标日前已发生的气象数据，预测次日参考蒸散，精度优于“沿用昨日值”这一业务上常用的简单基线；
\item \textbf{误差可量化}：为每个预测给出名义90\%的预测区间，并在未参与训练的年份上检验其实际覆盖率；
\item \textbf{预测进入决策}：将预测值或区间上界代入FAO-56根区水量平衡，计算是否灌溉、毛灌溉量、净补水量，并显式检查日供水能力；
\item \textbf{策略可比较}：在统一的气象与土壤情景下，比较定额日历、墒情阈值、昨日值预测、AI点预测与AI区间上界五种策略的用水量与缺水风险；
\item \textbf{过程可追溯}：交互界面、服务端计算与导出记录保持一致，全部结果能由原始数据和脚本重现。
\end{enumerate}

\subsection{技术路线与系统组成}
图\ref{fig:arch}给出系统的整体结构。系统分为数据层、预测层、决策层和交互层：数据层获取并校验三地公开气象序列；预测层构造49维历史特征，训练梯度提升树并在独立年份校准误差区间；决策层将参考蒸散乘以作物系数，计入根区亏缺，并在供水上限约束下计算补水量；交互层以图谱形式展示需水曲线、预测区间、根区水量和建议结果，并支持导出决策记录。

\begin{figure}[H]\centering
\begin{tikzpicture}[font=\small,
 box/.style={draw=forest,line width=.8pt,rounded corners=2pt,fill=white,align=center,minimum height=12mm,minimum width=27mm,inner xsep=3mm},
 layer/.style={fill=forest!8,draw=none,rounded corners=3pt},
 arr/.style={-{Stealth[length=2.4mm]},line width=.8pt,color=forest}]
\node[box] (d) {数据层\\[1pt]\footnotesize ERA5逐日气象\\\footnotesize 校验与归档};
\node[box,right=8mm of d] (m) {预测层\\[1pt]\footnotesize 49维特征\\\footnotesize 梯度提升树};
\node[box,right=8mm of m] (c) {不确定性\\[1pt]\footnotesize 2023年独立校准\\\footnotesize 90\%预测区间};
\node[box,right=8mm of c] (w) {决策层\\[1pt]\footnotesize FAO-56水量平衡\\\footnotesize 供水能力约束};
\node[box,below=10mm of m] (u) {交互层\\[1pt]\footnotesize 历史回放 · 参数调节 · 记录导出};
\node[box,below=10mm of w] (s) {策略仿真\\[1pt]\footnotesize 5种策略×36个情景};
\draw[arr] (d)--(m); \draw[arr] (m)--(c); \draw[arr] (c)--(w);
\draw[arr] (w.south)--(s.north);
\draw[arr] (w.south west)--(u.north east);
\draw[arr] (m.south)--(u.north);
\end{tikzpicture}
\caption{墒策系统结构与数据流}\label{fig:arch}
\end{figure}

\subsection{作品特色与创新点}
\textbf{（1）把预测误差放进灌溉决策。}多数需水预测工作只报告点预测精度。本作品用分割式共形方法在独立年份上校准误差分位数\cite{vovk,lei}，将区间上界作为“保守需水”直接代入灌溉触发规则，使使用者能够在“少用水”和“少缺水”之间作出知情选择，而不是被动接受一个数字。

\textbf{（2）机理模型与数据驱动方法耦合。}预测部分采用数据驱动的梯度提升树，决策部分严格遵循FAO-56的作物系数与根区水量平衡框架。模型只负责其擅长的短期气象需水预测，土壤、作物与设备约束由可解释的机理公式处理，二者通过作物系数与亏缺状态衔接，符合赛题倡导的“机理—数据耦合建模”。

\textbf{（3）供水能力显式建模。}当建议补水量超过设备或水源的日供水上限时，系统按上限给出可执行的补水量，同时提示能力不足与补水后的预计亏缺，而不是悄悄截断、营造“已满足需水”的假象。这一设计贴近灌区轮灌、机井出水量受限等实际约束。

\textbf{（4）严格的时间隔离与全链条可复现。}特征只使用目标日前的数据，训练、校准、测试按年份切分，并以自动化测试验证“扰动目标日及之后数据不改变输入特征”。原始接口响应、模型文件、逐日预测与180组仿真明细随工程交付，主要指标均可由脚本从CSV重算。

\textbf{（5）正负结果同时报告。}作品不仅给出整体指标，还按月份、地区和极端误差样本做了诊断，发现夏季高需水期区间覆盖不足等问题，并据此提出改进路线，使评估结论更可靠，也为后续优化指明了方向。
